% ============================================================
% ESTILO DE TÍTULOS — Variante A: título sobre fondo blanco
% El título deja de ser una "banda" de cabecera y pasa a leerse
% como contenido: texto oscuro en negrita sobre blanco, separado
% del cuerpo por un filete de color de acento.
%
% El conjunto (título + filete) se sitúa a la altura del bloque
% del logo de la barra lateral, igual que hacía la banda azul del
% tema, para no perder alto útil de diapositiva.
% Ajustar \alturaTitulo si se cambia el tamaño del logo.
% ============================================================
\definecolor{tituloTexto}{RGB}{20,33,97}    % azul UC3M oscuro
\definecolor{tituloAcento}{RGB}{214,118,0}  % ámbar de acento

\newlength{\alturaTitulo}
\setlength{\alturaTitulo}{-25pt}            % desplazamiento vertical del título

\setbeamercolor{frametitle}{fg=tituloTexto,bg=white}
\setbeamerfont{frametitle}{size=\large,series=\bfseries}

\setbeamertemplate{frametitle}{%
  \nointerlineskip%
  \vskip\alturaTitulo%
  \begin{beamercolorbox}[wd=\textwidth,leftskip=0.8em,rightskip=0.8em]{frametitle}
    \usebeamerfont{frametitle}\insertframetitle\strut\par
    \vskip3pt%
    {\color{tituloAcento}\rule{\dimexpr\textwidth-1.6em\relax}{1.4pt}}%
  \end{beamercolorbox}%
  \vskip6pt%
}
